\documentclass[tikz,border=4pt]{standalone}
\input{tikz-tcee}
\usepackage{pgfplotstable}
% colores-solucion.tex — Colores de los gráficos propios de la solución (1D01,
% examen de España). Se cargan después de tikz-tcee.tex (paleta «dic…»).
% Los 7 primeros son los que añade el cuadernillo del examen (mismo HEX, para
% que lo mismo tenga el mismo color en el examen y en la solución; CRITERIOS §7).
\definecolor{dicbienes}{HTML}{2C6E9B}
\definecolor{dicturismo}{HTML}{8CC3E6}
\definecolor{dicrentas}{HTML}{A6A6A6}
\definecolor{dicsaldo}{HTML}{1A1A1A}
\definecolor{dicpiin}{HTML}{8C510A}
\definecolor{dicfiscal}{HTML}{4A5A6A}
\definecolor{dicvivienda}{HTML}{B07AA1}
% Nuevos en la solución:
% sector privado (gama cálida, la del gasto privado de la guía)
\definecolor{dicprivado}{HTML}{C76A1F}
% efecto bola de nieve (azul medio, distinto del pizarra de lo fiscal)
\definecolor{dicbola}{HTML}{4F86C6}
% ajuste déficit-deuda y otras variaciones (residuos): gris claro con trama
\definecolor{dicresiduo}{HTML}{B4B4B4}
% previsión del PIB (el morado del PIB, aclarado)
\colorlet{dicpibprev}{dicpib!35!white}
\pgfplotsset{
  dicsol/.style={width=13cm, height=6.4cm, scale only axis=false,
    axis lines*=left, tick align=outside, clip=false,
    ymajorgrids, grid style={tceegris!25},
    ticklabel style={font=\footnotesize}, label style={font=\footnotesize},
    xticklabel style={/pgf/number format/1000 sep=},
    yticklabel style={/pgf/number format/.cd, use comma, 1000 sep={}},
    legend style={at={(0.5,-0.15)}, anchor=north, draw=none, font=\footnotesize,
      /tikz/every even column/.append style={column sep=8pt}},
    legend cell align=left},
  leyendalinea/.style={legend image code/.code={\draw[##1] (0cm,0cm) -- (0.55cm,0cm);}},
  % leyenda de una serie de solo puntos (en un gráfico con area legend)
  leyendamarca/.style={legend image code/.code={\draw[mark repeat=2,mark phase=2,##1] plot coordinates {(0cm,0cm) (0.3cm,0cm) (0.6cm,0cm)};}},
}

\begin{document}
% Solución, gráfico s08. Variación de la deuda pública (pp del PIB, puntos) y
% sus 3 componentes (barras): déficit primario, efecto bola de nieve
% (i - g)/(1 + g)·d(t-1) y ajuste déficit-deuda (residuo), 2021-2027
% (2026-2027: previsión de la Comisión). Datos: ../datos/s08-deuda.csv (cuadro 9).
% Capas: sombra de la previsión → barras → cero → puntos → rótulos.
\pgfplotstableread[col sep=comma,comment chars={\#}]{../datos/s08-deuda.csv}\deuda
\begin{tikzpicture}[tcee]
  \begin{axis}[dicsol, width=11.5cm, height=6.4cm,
      ybar stacked, stack negative=separate, bar width=0.6cm,
      xmin=2020.4, xmax=2027.6, xtick={2021,...,2027},
      ymin=-10, ymax=6, ytick={-10,-8,...,6},
      ylabel={puntos del PIB},
      legend columns=2, area legend, dicapas]
    \capa{1}{\dicsombra[fill=tceegris!12]{2025.5}{-10}{2027.6}{6}
      \dicrotulo[font=\scriptsize\itshape, anchor=north]{(2026.5,5.8)}{Previsión}}
    \capa{2}{\addplot[on layer=dicbarras, dicbarra=dicfiscal] table[x=anio, y=primario] {\deuda};
      \addlegendentry{Déficit primario (+) o superávit ($-$)}}
    \capa{3}{\addplot[on layer=dicbarras, dicbarra=dicbola] table[x=anio, y=bola] {\deuda};
      \addlegendentry{Bola de nieve: $\frac{i-\gamma^N}{1+\gamma^N}\,d_{t-1}$}}
    \capa{4}{\addplot[on layer=dicbarras, dicbarratrama=dicresiduo] table[x=anio, y=add] {\deuda};
      \addlegendentry{Ajuste déficit-deuda (residuo)}}
    \capa{5}{\dicraya{(2020.4,0) (2027.6,0)}}
    \capa{6}{\addplot[on layer=dicpuntos, stack plots=false, only marks, mark=diamond*, mark size=3.3pt,
        mark options={fill=dicsaldo, draw=white, line width=0.4pt}, leyendamarca] table[x=anio, y=var_deuda] {\deuda};
      \addlegendentry{Variación de la deuda ($\Delta d$)}}
    \capa{7}{
      \dicrotulo[font=\scriptsize, anchor=west, inner sep=1pt]{(2021.27,-3.6)}{$-$3,6}
      \dicrotulo[font=\scriptsize, anchor=west, inner sep=1pt]{(2022.27,-6.4)}{$-$6,4}
      \dicrotulo[font=\scriptsize, anchor=west, inner sep=1pt]{(2023.27,-4.0)}{$-$4,0}
      \dicrotulo[font=\scriptsize, anchor=west, inner sep=1pt]{(2024.27,-3.9)}{$-$3,9}
      \dicrotulo[font=\scriptsize, anchor=west, inner sep=1pt]{(2025.27,-0.9)}{$-$0,9}
      \dicrotulo[font=\scriptsize, anchor=west, inner sep=1pt]{(2026.27,-0.9)}{$-$0,9}
      \dicrotulo[font=\scriptsize, anchor=west, inner sep=1pt]{(2027.27,-0.7)}{$-$0,7}
    }
  \end{axis}
\end{tikzpicture}
\end{document}
